\uuid{7K8i}
\titre{Tirage au sort de SMS}
\auteur{worms}
\organisation{exo7}
\licence{CC BY-NC-SA 4.0}
\datecreate{2017-08-28}
\niveau{L2}
\module{Probabilité et statistique}
\chapitre{Probabilité discrète}
\sousChapitre{Probabilité et dénombrement}
\difficulte{}
\theme{Probabilités, événements}
\exo7qcmid{worms.worms.010}
\serie{ordonnee}

\contenu{
\texte{
On suppose que 2000 personnes ont envoyé un SMS dans le cadre d'un mini-jeu télé qui consistait à répondre à une question (particulièrement idiote) à 2 choix. On suppose que la société qui gère ce "jeu-SMS" sélectionne 30 SMS au hasard parmi les 2000.
}

\begin{question}
\qclassification{260.01}
Dans cette question et les suivantes, on note $\Omega$ l'ensemble de tous les sous-ensembles de 30 SMS (distincts). Combien d'éléments $\Omega$ contient-il ?
\begin{answers}
\good{$C_{2000}^{30}$}
\bad{$A_{2000}^{30}$}
\bad{1971}
\bad{$2000!/30!$}
\end{answers}
\begin{explanations}
C'est $C_{2000}^{30}$, voir le cours !
\end{explanations}
\end{question}

\begin{question}
\qclassification{260.01}
A-t-on équiprobabilité dans cette situation ?
\begin{answers}
\good{Oui.}
\bad{Non.}
\end{answers}
\begin{explanations}
Bien sûr que oui. Tous les sous-ensembles de 30 SMS ont autant de chances d'être tirés.
\end{explanations}
\end{question}

\begin{question}
\qclassification{260.01}
On considère maintenant que vous faites partie des personnes qui ont envoyé un SMS. Quelle est la probabilité que vous soyez sélectionné(e) ?
\begin{answers}
\good{$1.5\%$}
\bad{$0.03\%$}
\bad{$0.15\%$}
\bad{Environ $1.2\%$.}
\end{answers}
\begin{explanations}
L'événement en question correspond au sous-ensemble de $\Omega$ contenant tous les sous-ensembles de 30 SMS dont le sien, et cet événement est de cardinal $C_{1999}^{29}$, et cela donne, après simplifications, $\mathbb{P}(A) = 30/2000 = 1,5\%$.
\end{explanations}
\end{question}

\begin{question}
\qclassification{260.01}
\qkeeporder
Votre ami(e) fait également partie des personnes ayant envoyé un SMS. Quelle est la probabilité qu'au moins l'un(e) d'entre vous soit tiré(e) au sort ?
\begin{answers}
\bad{Deux fois la réponse à la question précédente.}
\good{Environ 3\% (mais pas 3\%).}
\bad{Environ 0,3\% (mais pas 0,3\%).}
\bad{Une autre valeur.}
\end{answers}
\begin{explanations}
Le cardinal du complémentaire $B^c$ de l'événement étudié est clairement $C_{1998}^{30}$,  nombre de sous-ensembles de 30 éléments parmi les 1998 SMS restants. Cela donne $\mathbb{P}(B)=1- (1970\times 1969)/(2000\times 1999) \simeq  2.98\%$.
\end{explanations}
\end{question}
}
